\subsection{Norms and traces in extensions of finite degree}

The transitivity formulae in algebras (III, p.~548) imply the following proposition in the case of extensions of finite degree.

PROPOSITION 2. --- Let F be an extension of finite degree of K and E a subextension of F. Then for every \(x \in F\) we have

\[
\operatorname{Tr} _ {\mathrm{F} / \mathrm{K}} (x) = \operatorname{Tr} _ {\mathrm{E} / \mathrm{K}} (\operatorname{Tr} _ {\mathrm{F} / \mathrm{E}} (x))\tag{6}
\]

\[
\mathrm{N} _ {\mathrm{F} / \mathrm{K}} (x) = \mathrm{N} _ {\mathrm{E} / \mathrm{K}} (\mathrm{N} _ {\mathrm{F} / \mathrm{E}} (x))\tag{7}
\]

COROLLARY. --- Put \(m = [F : E]\) ; then for every \(x \in E\) we have

\[
\operatorname{Tr} _ {\mathrm{F} / \mathrm{K}} (x) = m. \operatorname{Tr} _ {\mathrm{E} / \mathrm{K}} (x)\tag{8}
\]

\[
\mathrm{N} _ {\mathrm{F} / \mathrm{K}} (x) = \mathrm{N} _ {\mathrm{E} / \mathrm{K}} (x) ^ {m}\tag{9}
\]

PROPOSITION 3. --- Let E be an extension of finite degree n of K and x an element of E, of degree d over K. Write \(f(X) = X^d + \sum_{i=1}^{d} a_i X^{d-i}\) for the minimal polynomial of x over K. Then we have

\[
\mathrm{Tr} _ {\mathrm{E/K}} (x) = - \frac {n}{d} a _ {1}\tag{10}
\]

\[
\mathrm{N} _ {\mathrm{E} / \mathrm{K}} (x) = ((- 1) ^ {d} a _ {d}) ^ {n / d} = (- 1) ^ {n} a _ {d} ^ {n / d}\tag{11}
\]

Prop. 3 follows directly from the Cor. of Prop. 2 and the lemma :

Lemma 3. --- Let R be a commutative ring, \(f(\mathbf{X}) = \mathbf{X}^d + \sum_{i=1}^{d} a_i \mathbf{X}^{d-i}\) a monic polynomial of R{[}X{]}, A the R-algebra R{[}X{]}/(f) and x the residue class of X in A. Then \(\operatorname{Tr}_{\mathrm{A/R}}(x) = -a_1\) and \(\mathrm{N}_{\mathrm{A/R}}(x) = (-1)^d a_d\) .

By the Cor. (IV, p.~11) the sequence \((1, x, \ldots, x^{d-1})\) is a basis of \(\mathbf{A}\) ; further we have

\[
\begin{array}{r l} x \cdot 1 = x, & x \cdot x = x ^ {2}, \dots , x \cdot x ^ {d - 2} = x ^ {d - 1}, \\ & x \cdot x ^ {d - 1} = - a _ {d} \cdot 1 - a _ {d - 1} \cdot x - \dots - a _ {1} \cdot x ^ {d - 1}. \end{array}
\]

The matrix which expresses the multiplication by x with respect to the basis \((1, x, \ldots, x^{d-1})\) of A is thus of the following form (we have taken d = 5 to fix the ideas):

\[
\left( \begin{array}{c c c c c} 0 & 0 & 0 & 0 & - a _ {5} \\ 1 & 0 & 0 & 0 & - a _ {4} \\ 0 & 1 & 0 & 0 & - a _ {3} \\ 0 & 0 & 1 & 0 & - a _ {2} \\ 0 & 0 & 0 & 1 & - a _ {1} \end{array} \right).
\]

The trace of this matrix is clearly \(-a_{1}\) ; the determinant may be calculated by expanding by the first row, and we then find

\[
(- 1) ^ {d - 1} (- a _ {d}) = (- 1) ^ {d} a _ {d}.
\]

For the rest of this No.~we denote by E an extension of finite degree of K and by x an element of E. We shall indicate how to calculate the norm and trace of x in various cases.

\begin{enumerate}
\def\labelenumi{\alph{enumi})}
\tightlist
\item
  Case of a separable extension: suppose that E is separable of degree n over K, denote by \(\Omega\) an algebraic closure of K and by \(\sigma_{1}, \ldots, \sigma_{n}\) the n distinct K-homomorphisms of E into \(\Omega\) . By Formula (3) (V, p.~48) we have in \(\Omega\)
\end{enumerate}

\[
\mathrm{Tr} _ {\mathrm{E} / \mathrm{K}} (x) = \sum_ {i = 1} ^ {n} \sigma_ {i} (x), \quad \mathrm{N} _ {\mathrm{E} / \mathrm{K}} (x) = \prod_ {i = 1} ^ {n} \sigma_ {i} (x).\tag{12}
\]

\begin{enumerate}
\def\labelenumi{\alph{enumi})}
\setcounter{enumi}{1}
\tightlist
\item
  Case of a p-radical extension: suppose that K is of characteristic p \textgreater{} 0 and that the extension E is p-radical; there exists an integer \(e \geqslant 0\) such that \([E : K] = p^{e}\) (V, p.~26, Prop. 4). If f is the height of x over K, the minimal polynomial of x over K is \(X^{p^{f}} - x^{p^{f}}\) (V, p.~24, Prop. 1). By Prop. 3 we have \(\mathrm{N}_{\mathrm{E}/\mathrm{K}}(x) = (x^{p^{f}})^{p^{e}/p^{f}}\) , whence
\end{enumerate}

\[
\mathrm{N} _ {\mathrm{E} / \mathrm{K}} (x) = x ^ {p ^ {e}} = x ^ {[ \mathrm{E}: \mathrm{K} ]}.\tag{13}
\]

For the trace we find \(\mathrm{Tr}_{\mathrm{E} / \mathrm{K}}(x) = -p^{e - f}a\) , where \(a\) is the coefficient of \(\mathbf{X}^{p^f - 1}\) in the polynomial \(\mathbf{X}^{p^f} - x^{p^f}\) ; in other words, we have

\[
\operatorname{Tr} _ {\mathrm{E} / \mathrm{K}} (x) = p ^ {e} \cdot x = [ \mathrm{E}: \mathrm{K} ] x = \left\{ \begin{array}{l l} x & \text { if } \quad [ \mathrm{E}: \mathrm{K} ] = 1 \\ 0 & \text { if } \quad [ \mathrm{E}: \mathrm{K} ] > 1. \end{array} \right.\tag{14}
\]

\begin{enumerate}
\def\labelenumi{\alph{enumi})}
\setcounter{enumi}{2}
\tightlist
\item
  General case: we can summarize the calculation of norm and trace in the following proposition:
\end{enumerate}

PROPOSITION 4. --- Let p be the characteristic exponent of K and E an extension of finite degree of K. Let \(\sigma_{1}, \ldots, \sigma_{n}\) be the distinct K-homomorphisms of E into an algebraic closure \(\Omega\) of K, and let \(p^{e} = [E : K]_{i}\) . For each \(x \in E\) we have in \(\Omega\)

\[
\operatorname{Tr} _ {\mathrm{E} / \mathrm{K}} (x) = p ^ {e} \cdot \sum_ {i = 1} ^ {n} \sigma_ {i} (x), \quad \mathrm{N} _ {\mathrm{E} / \mathrm{K}} (x) = \left(\prod_ {i = 1} ^ {n} \sigma_ {i} (x)\right) ^ {p ^ {e}}.\tag{15}
\]

Let \(E_s\) be the relative separable closure of K in E; then \(E_s\) is a separable extension of degree n of K and \(\sigma_1, \ldots, \sigma_n\) induce distinct K-homomorphisms of \(E_s\) into \(\Omega\) ; further, E is a p-radical extension of \(E_s\) of degree \(p^e\) (V, p.~44, Prop. 13 and p.~46). So Prop. 4 follows from the Formulae (6), (7), (13), (14) and (12).

